\documentclass[10pt,a4paper]{amsart}
\usepackage[margin=19mm]{geometry}
\usepackage{amsmath,amssymb,mathtools}
\usepackage[scr=rsfs,cal=euler]{mathalfa}
\usepackage{xfrac}
\usepackage[unicode,hidelinks,bookmarks=false]{hyperref}
\newcommand{\ie}{i.e.\ }
\newcommand{\cf}{cf.\ }
\newcommand{\resp}{respectively\ }
\newcommand{\Subsection}{\subsection}
\providecommand{\nobreakdash}{\nobreak\hbox{-}}
\setlength{\emergencystretch}{2em}
\multlinegap0pt
\usepackage{tikz}
\usetikzlibrary{cd}
\usepackage{xcolor}
\usepackage{amssymb}
\usepackage[normalem]{ulem}
\usepackage[shortlabels]{enumitem}
\usepackage{float}
\usepackage{dsfont}
\usepackage{csquotes}
\usepackage{algorithm}\usepackage{algpseudocode}
\usepackage{mathtools}
\newtheorem{thm}{Theorem}
\title{Selfless inclusions arising from commensurator groups of hyperbolic groups}
\author{Aaratrick Basu, Felipe Flores}
\date{Source: arXiv:2605.12737v1 (CC BY 4.0)}
\begin{document}
\maketitle

\noindent\emph{Source and licence.} Selfless inclusions arising from commensurator groups of hyperbolic groups. Version arXiv:2605.12737v1 (CC BY 4.0), distributed by arXiv under the Creative Commons Attribution 4.0 International licence, http://creativecommons.org/licenses/by/4.0/, as recorded in the arXiv licence metadata for this exact version and shown on the abstract page https://arxiv.org/abs/2605.12737v1. Full licence notice: \texttt{licenses/arxiv-2605.12737-CC-BY-4.0.txt}.\par\smallskip

\noindent\textit{Editorial scope.} Counted unit: the article's thm thm (source0.tex lines 441--443). The article's complete original proof of it is delivered with it (source0.tex lines 444--472, one \texttt{proof} environment of 29 lines). No mathematics is written, repaired or re-derived: the statement and the proof are the article's own lines, in the article's own notation, and no retained line is substituted or re-flowed.  No further source-local context is needed: every cross-reference of every retained line resolves inside the two delivered spans.  Nothing was omitted from any retained span, so the package carries no omission record.  No retained line contains a citation, so the wrapper carries no bibliography and no bibliography entry.  The retained lines contain the article's own diagram code, which the wrapper typesets with the standard tikz packages; no image file is delivered and the package declares no asset dependency.  Editorial formatting only, in addition: an AMS \texttt{amsart} preamble that restates the article's own declarations and macros used by the retained lines, namely the environments \texttt{thm} on separate counters, together with the standard packages those lines need.  The retained environments print in this document as Theorem 1 in order of appearance, whereas the article prints the same items with the numbering of its own full document.  The complete native source of the article as distributed in the arXiv e-print for this exact version is bundled unchanged as \texttt{sources/200441/source0.tex}.
\begin{thm}
   Let \( H \leq_c G \) be a commensurated inclusion, where \( H \) is hyperbolic. The definition above defines an action of \( G \) on \( \partial H \) that restricts to that of \( H \) on \( \partial H \) induced by conjugation on itself.
\end{thm}

\begin{proof}
    We get \( g_* \) is the same homeomorphism as that induced by conjugation by \( g \) whenever \( g \in H \), since the relevant inclusion maps are indentities. We now check that the definition is indeed an action of \( G \), i.e, \( (g_1 g_2)_* = (g_1)_* (g_2)_* \) for all \( g_1, g_2 \in G \). The essential observation is that 
   \[
       K = H \cap g_1^{-1}Hg_1 \cap g_2^{-1}Hg_2 \cap g_2^{-1}g_1^{-1}Hg_1g_2
   \]
   is an intersection of finite index subgroups of \( H \), and thus is still of finite index, along with the fact that \( \iota_{g_1}, \iota_{g_2}, \iota_{g_1g_2} \) satisfy \( \iota_{g_1g_2} = \iota_{g_1} \iota_{g_2} \), and in particular they induce the same map \( \partial \iota \) on boundaries. We introduce \( H_i = H \cap g^{-1}_i H g_i, i = 1,2, \) and \( H_{12} = H \cap (g_1g_2)^{-1} H g_1g_2 \), so that \( K = H_1 \cap H_2 \cap H_{12} \). We then have a diagram:
\[\begin{tikzcd}
	&& \partial H_{12} && \partial \bar H_{12} && \\
	\partial H && \partial K && \partial \bar K && \partial H \\
	& \partial H_2 & \partial \bar H_2 & \partial H & \partial H_1 & \partial \bar H_1
	\arrow[from=1-3, to=1-5, "\partial \iota_{g_1 g_2}"]
	\arrow[from=1-5, to=2-7]
	\arrow[from=2-1, to=1-3]
	\arrow[from=2-1, to=2-3]
	\arrow[from=2-1, to=3-2]
	\arrow[from=2-3, to=1-3]
	\arrow[from=2-3, to=2-5, "\partial \iota"]
	\arrow[from=2-3, to=3-2]
	\arrow[from=2-5, to=1-5]
	\arrow[from=2-5, to=2-7]
	\arrow[from=2-5, to=3-6]
	\arrow[from=3-2, to=3-3, "\partial \iota_{g_2}"']
	\arrow[from=3-3, to=3-4]
	\arrow[from=3-4, to=3-5]
	\arrow[from=3-5, to=3-6, "\partial \iota_{g_1}"']
	\arrow[from=3-6, to=2-7]
\end{tikzcd}\]
Here the bar on top indicates the image group under the respective conjugations, and the unlabelled maps are the boundary maps induced by inclusions or their inverses. By functoriality of maps induced on boundaries, the way inclusions compose, and the fact that all of the maps in the diagram are homeomorphisms, we get the diagram is commutative, and thus \( (g_1g_2)_* = (g_1)_* (g_2)_* \) as was required.
\end{proof}

\end{document}
